\subsection{P018 --- Semantic model-based large-scale deployment of AI-driven building management applications}
\label{paper:P018}
\textbf{Authors:} Kan Xu, Zhe Chen, Fu Xiao, Jing Zhang, Hanbei Zhang, Tianyou Ma\par
\textbf{Major Category:} BIM and Semantic Information\par
\textbf{Secondary Category:} AI and Machine Learning for Buildings, BMS, BAS, BACS and Building Automation\par
\textbf{Keywords:} Semantic model, Brick Schema, AI applications, Building management, Data interoperability, Large-scale deployment\par
\paragraph{主な主張}
Brickに基づくsemantic modelでBIM, BMS, IoT由来のstatic dataとtemporal dataを統合し, SPARQLとFLUXによる標準化されたdata acquisitionを介して同一application workflowを3建物へ展開した. 24 h ahead cooling-load predictionのCV-RMSEはSite1 14.5\%, Site2 19.5\%, Site3 13.5\%であった. \cite{Xu:2024SemanticAI}
\paragraph{新規性}
既往のsemantic model-based building applicationsが単一建物またはopen dataset中心であったのに対し, graph databaseとtemporal databaseをsemantic identifierで接続するframeworkを構築し, educational, office, commercialの3実建物でdeploymentを実証した点に新規性がある.
\paragraph{位置付け}
semantic informationを予測モデルへの追加featureとして扱う研究ではなく, heterogeneous building dataをbuilding-independentに取得するabstraction layerとして利用し, AI applicationのportabilityとlarge-scale deploymentを支援する研究として位置付ける.
